% !TEX program = xelatex
\documentclass[12pt,a4paper]{ctexart}

% ========== 消除 Fandol 字体警告 ==========
\PassOptionsToPackage{Script=CJK}{fontspec}

% ========== 宏包 ==========
\usepackage[margin=2.5cm]{geometry}
\usepackage{hyperref}
\hypersetup{colorlinks=true,linkcolor=blue,urlcolor=teal}
\usepackage{xcolor}
\usepackage{listings}
\usepackage{tcolorbox}
\tcbuselibrary{breakable,skins}
\usepackage{fontawesome5}

% ========== 代码样式 ==========
\lstset{
  basicstyle=\ttfamily\small,
  breaklines=true,
  backgroundcolor=\color{gray!10},
  frame=single,
  rulecolor=\color{gray},
  columns=flexible,
  xleftmargin=2em,
  xrightmargin=2em,
}

% ========== 自定义提示框 ==========
\newtcolorbox{notebox}[1][]{
  enhanced, breakable,
  colback=blue!5!white, colframe=blue!75!black,
  fonttitle=\bfseries,
  title=#1,
  attach boxed title to top left={xshift=2mm,yshift=-2mm},
  boxed title style={colback=blue!75!black,size=small},
  before skip=12pt plus 2pt,
  after skip=12pt plus 2pt,
}

\newtcolorbox{warnbox}[1][]{
  enhanced, breakable,
  colback=red!5!white, colframe=red!75!black,
  fonttitle=\bfseries,
  title=#1,
  attach boxed title to top left={xshift=2mm,yshift=-2mm},
  boxed title style={colback=red!75!black,size=small},
  before skip=12pt plus 2pt,
  after skip=12pt plus 2pt,
}

% ========== 文档信息 ==========
\title{\textbf{LaTeX 编译与中文字体疑难解析}\\{\large ——计算机专业学习笔记}}
\author{你的名字}
\date{\today}

\begin{document}
\maketitle
\tableofcontents
\newpage

% -------------------------------------------------
\section{从源码到 PDF：LaTeX 的“翻译”流程}
% -------------------------------------------------

我们可以把一份 \texttt{.tex} 源码想象成一封用\emph{世界通用语}写的信，
而 \LaTeX 引擎则是一位\emph{翻译官}。信件内容完全一样，但交给不同的翻译官，
他们会用不同的方式表达出来，最终生成一份精美的 PDF 文件。

基本流程如下：
\begin{enumerate}
  \item 用任意文本编辑器（vim、nano、VS Code 等）写一个 \texttt{.tex} 文件。
  \item 在终端调用某个 \LaTeX 引擎进行编译。
  \item 引擎读入 \texttt{.tex} 源码，处理其中的命令和环境，输出 \texttt{.pdf}。
\end{enumerate}

\begin{notebox}[最简单的英文文档]
\begin{lstlisting}[language={[LaTeX]TeX}]
\documentclass{article}
\begin{document}
Hello, \LaTeX!
\end{document}
\end{lstlisting}
编译命令：\texttt{pdflatex hello.tex}
\end{notebox}

\begin{notebox}[支持中文的文档]
\begin{lstlisting}[language={[LaTeX]TeX}]
\documentclass{ctexart}   % 或者 \documentclass[UTF8]{ctexart}
\begin{document}
你好，\LaTeX！计算机学习快乐！
\end{document}
\end{lstlisting}
编译命令：\texttt{xelatex hello.tex}
\end{notebox}

\section{认识三位翻译官：pdfTeX、XeTeX、LuaTeX}
\label{sec:engines}

这三位“翻译官”各有所长，理解它们的区别是举一反三的关键。

\begin{description}
  \item[pdfTeX（老牌翻译官）] \hfill \\
  最传统的引擎，直接输出 PDF。它最擅长处理\textbf{英文}和传统 TeX 字体，
  对 Unicode 的支持非常有限，处理中文时就像一个只会说方言的老翻译官——需要借助
  CJK 宏包和各种专门的字体包，稍有不慎就会“翻译失败”。

  \item[XeTeX（现代多语种翻译官）] \hfill \\
  一位精通多国语言的现代翻译官。它原生支持 UTF-8 编码，能直接调用你操作系统里
  安装的任何字体（如宋体、黑体、Noto 等）。处理中文时，XeTeX 是\textbf{首选}。

  \item[LuaTeX（全能编程式翻译官）] \hfill \\
  XeTeX 的“升级版”，同样支持 Unicode 和系统字体，而且内部嵌入了一个 Lua 脚本
  解释器，允许用编程方式深度定制排版逻辑。功能非常强大，但日常使用中，XeTeX
  已经足够。
\end{description}

\begin{warnbox}[关键结论]
  简单来说：
  \begin{itemize}
    \item 纯英文文档、老式论文模板：用 \texttt{pdflatex}
    \item 带中文、多语种的文档：用 \texttt{xelatex}（推荐）或 \texttt{lualatex}
  \end{itemize}
\end{warnbox}

\section{中文排版为何首选 XeLaTeX？}
\label{sec:why-xetex}

这就要讲到字体处理的\textbf{编码模型}了。

\begin{itemize}
  \item \textbf{pdfTeX 的世界}：传统 TeX 字体每个字符占 1 个字节（最多 256 个字符位置）。
        英文绰绰有余，但中文有上万个汉字，必须把字体分割成多个小字库，还要通过
        \texttt{.tfm} 文件和映射关系来定位。一旦某个字体包没装好，就会出“fandol 不可用”
        这类错误。
  \item \textbf{XeTeX 的世界}：直接使用 \textbf{UTF-8 编码}，每个字符有全局唯一的编码，
        字体的选择直接交给操作系统。只要你的系统里安装了中文字体，XeTeX 就能正确显示。
        这就好比打电话：pdfTeX 需要提前查内部黄页，而 XeTeX 直接拨号。
\end{itemize}

因此在实践中，\texttt{ctex} 宏包搭配 XeTeX 可以省去无数字体配置的烦恼。

\section{实战排错：fandol 字体缺失错误}
\label{sec:fandol-error}

你是不是在终端敲了 \texttt{pdflatex fedora-install-crewai-note.tex}，然后看到了
下面这样的错误？

\begin{lstlisting}[breaklines]
! Critical Class ctexart Error: CTeX fontset `fandol' is unavailable in
(ctexart)                       current mode.
\end{lstlisting}

\textbf{错误翻译}：“ctexart 文档类严重错误：在当前模式下，CTeX 字体集 \texttt{fandol}
不可用。”

\subsection{为什么会这样？}

因为你在用 \texttt{pdflatex}（老翻译官）编译一个使用了 \texttt{ctexart} 的中文文档。
\texttt{ctexart} 在 pdfTeX 模式下默认选择 \textbf{fandol} 字体集，它需要安装对应的
\texttt{texlive-fandol} 包，但你的 Fedora 系统上并没有装这个包。

更本质的原因是：\texttt{fandol} 是为 pdfTeX 准备的一套 Type1 中文字体，而很多 TeX
发行版的最小安装中并不包含它。

\subsection{两种解决方案}
\label{sec:solutions}

\begin{notebox}[方案一：换个翻译官（推荐）]
直接改用现代引擎 \texttt{xelatex}：
\begin{lstlisting}[language=bash]
xelatex fedora-install-crewai-note.tex
\end{lstlisting}
优点：
\begin{itemize}
  \item 彻底避免 pdfTeX 字体问题
  \item 直接利用系统字体，兼容各种语言文字
  \item 编译速度通常也更快（中文部分）
\end{itemize}
\end{notebox}

\begin{warnbox}[方案二：给老翻译官配字典（备选）]
安装缺失的 fandol 字体包：
\begin{lstlisting}[language=bash]
sudo dnf install texlive-fandol
\end{lstlisting}
然后再用 \texttt{pdflatex} 编译。如果想用其他字体集，也可以修改文档类选项：
\begin{lstlisting}[language={[LaTeX]TeX}]
\documentclass[fontset=ubuntu]{ctexart}  % 使用文泉驿字体，需提前安装
\end{lstlisting}
\end{warnbox}

\subsection{我是怎么诊断出来的？}
\label{sec:debug}

\begin{enumerate}
  \item \textbf{看日志}：\texttt{.log} 文件里明明白白写着 \texttt{fandol} 不可用。
  \item \textbf{查引擎}：确认一下自己是不是误用了 \texttt{pdflatex}。
  \item \textbf{试切换}：有条件的优先换 \texttt{xelatex}，一般立刻解决问题。
  \item \textbf{搜包名}：实在要用 \texttt{pdflatex} 时，用 \texttt{dnf search} 搜索
        \texttt{fandol} 或 \texttt{texlive-lang-chinese} 来补装字体。
\end{enumerate}

\section{ctex 宏包的字体集机制}
\label{sec:fontset}

\texttt{ctex} 文档类（\texttt{ctexart, ctexbook, ctexrep}）会根据\textbf{编译引擎}自动
选择一套中文字体集。这套机制就像一个聪明的管家：

\begin{center}
\begin{tabular}{p{2.5cm} p{4cm} p{4cm}}
  \hline
  \textbf{引擎} & \textbf{默认字体集} & \textbf{背后的字体调用} \\
  \hline
  pdfLaTeX & fandol (未安装则报错) & CJK 宏包 + fandol Type1 字体 \\
  XeLaTeX  & 系统字体（如宋体、黑体） & fontspec 宏包 + 操作系统字体 \\
  LuaLaTeX & 系统字体 & fontspec 宏包 + 操作系统字体 \\
  \hline
\end{tabular}
\end{center}

你可以用 \texttt{fontset} 选项手动指定字体集，例如：
\begin{lstlisting}[language={[LaTeX]TeX}]
\documentclass[fontset=windows]{ctexart}   % 需要 Windows 系统字体
\documentclass[fontset=mac]{ctexart}       % 需要 macOS 系统字体
\documentclass[fontset=fandol]{ctexart}    % 强制使用 fandol
\end{lstlisting}

在 Fedora 等 Linux 系统下，最省心的做法就是\textbf{不指定 fontset，直接用 xelatex}，
让管家自动调用系统中已安装的文泉驿、Noto 等开源字体。

\section{另一个常见陷阱：\texttt{\textbackslash end} 不能在 \texttt{\textbackslash texttt} 里裸用}
\label{sec:end-in-texttt}

\subsection{事故现场}
在编写这份笔记时，我们遇到了一个奇怪的编译错误，按回车后还无限循环：
\begin{lstlisting}[breaklines]
! Missing \endcsname inserted.
<to be read again>
                   \aftergroup
l.270 ...nning use of ...} → 缺少右花括号或 \texttt{\end}
\end{lstlisting}

出错的这一行想表达的意思其实是：“缺少右花括号或 \texttt{\textbackslash end}”，但在代码里我们直接写了：
\begin{lstlisting}[language={[LaTeX]TeX}]
... 缺少右花括号或 \texttt{\end}
\end{lstlisting}

\subsection{为什么会爆炸？}
\label{sec:why-end-explodes}

在 LaTeX 中，\verb|\end| 是一个\emph{控制序列}，它的作用是结束当前环境（如 \verb|\end{document}|、\verb|\end{itemize}|）。当你把它放进 \verb|\texttt{...}| 这类命令的参数里时，LaTeX 会先尝试\emph{展开}（执行）这个控制序列，而不是把它当成普通文字打印出来。因为此时并没有一个合理的环境可以结束，内部机制就会彻底混乱，抛出 \texttt{Missing \textbackslash endcsname} 错误。

简单来说：\textbf{控制序列 \texttt{\textbackslash end} 在参数里“越狱”了，破坏了 LaTeX 的正常语法结构。}

\subsection{安全打印控制序列的方法}
\label{sec:safe-print}

\begin{enumerate}
  \item \textbf{用 \texttt{\textbackslash verb}} （最推荐）\\
  \verb|\verb| 会原样输出其中的所有字符，不解释任何命令。缺点是它不能出现在其他命令的参数中，所以通常需要单独使用。例如：
  \begin{lstlisting}[language={[LaTeX]TeX}]
  ... 缺少右花括号或 \verb|\end|，找一个花括号匹配插件...
  \end{lstlisting}

  \item \textbf{用 \texttt{\textbackslash texttt\{\textbackslash textbackslash end\}}}\\
  先用 \verb|\textbackslash| 输出一个反斜杠符号，再跟上普通文字 \texttt{end}，这样就不会触发控制序列了。

  \item \textbf{用 \texttt{\textbackslash string}}\\
  \verb|\string| 会让紧随其后的控制序列丧失功能，变成普通字符序列。例如：
  \begin{lstlisting}[language={[LaTeX]TeX}]
  \texttt{\string\end}
  \end{lstlisting}
\end{enumerate}

\subsection{举一反三}
\label{sec:general-special-chars}

这个原则不仅适用于 \verb|\end|，对所有 LaTeX 特殊字符（如 \verb|\begin|、\verb|{|、\verb|}|、\verb|%|、\verb|#| 等）都成立。当你想在普通文字（或等宽字体）中显示它们时，必须使用对应的转义命令或 \verb|\verb|。养成习惯后，这类“诡异”的编译错误一眼就能看穿。

\section{搭建高效的 LaTeX 编译工作流}
\label{sec:workflow}

\subsection{工具选择}
\begin{itemize}
  \item \textbf{纯命令行}：vim/emacs + 终端，适合极客
  \item \textbf{现代编辑器}：VS Code 安装 LaTeX Workshop 插件，自动编译和预览
  \item \textbf{专用 IDE}：TeXstudio、TeXmaker，Fedora 下 \texttt{sudo dnf install texstudio}
\end{itemize}

\subsection{编译自动化}
如果文档有目录、交叉引用或参考文献，需要多次编译。可以用自动化工具：
\begin{itemize}
  \item \texttt{latexmk}：智能判断需要编译几次，例如\\
        \texttt{latexmk -pdf -xelatex mydoc.tex}\\
        它会自动调用 \texttt{xelatex} 并运行必要的 \texttt{bibtex} 等。
  \item 编写 Makefile 或脚本，适合大项目。
\end{itemize}

\section{举一反三：常见编译错误与调试心法}
\label{sec:general-debug}

掌握了上面这些，你就能解决一大类“找不到字体/宏包”的问题。这里再总结一个通用排错框架。

\begin{notebox}[遇到错误不要慌，三问自己]
\begin{enumerate}
  \item \textbf{引擎用对了吗？} \\
  中文文档用 \texttt{xelatex}，带特定字体要求的模板可能只能用 \texttt{pdflatex}。
  \item \textbf{宏包/字体装了吗？} \\
  类似 \texttt{! LaTeX Error: File `xxx.sty' not found.} 就是宏包缺失；
  \texttt{fandol} 不可用则是字体缺失。Fedora 下用 \texttt{dnf search} 搜索安装。
  \item \textbf{源码本身有语法错误吗？} \\
  少个花括号、命令拼写错误、数学模式内外混乱都会引发连锁报错。顺着第一个错误往上找。
\end{enumerate}
\end{notebox}

\subsection{常见错误速查}
\begin{itemize}
  \item \texttt{Undefined control sequence} → 命令拼写错误或宏包未加载。
  \item \texttt{Misplaced alignment tab character \&} → 表格/矩阵的 \texttt{\&} 写在了外面，检查是否漏了环境。
  \item \texttt{File ended while scanning use of ...} → 缺少右花括号或 \verb|\end|，找一个花括号匹配插件可快速定位。
  \item \texttt{Font ... not found} → XeLaTeX 下缺少系统字体，安装对应字体或换字体集。
\end{itemize}

\subsection{一个更直观的类比}
把 LaTeX 编译想象成\textbf{盖房子}：
\begin{itemize}
  \item \texttt{.tex} 源文件是\textbf{建筑设计图}。
  \item 各种宏包和字体是\textbf{建筑材料}。
  \item 编译引擎是\textbf{施工队}（pdfTeX 施工队、XeTeX 施工队）。
  \item 报错就是施工过程中发现图纸矛盾，或材料缺货。
\end{itemize}
图纸没问题、材料齐全、找对施工队，房子（PDF）自然就盖好了。

\section{结语}
至此，你不仅知道了如何创建和编译 LaTeX 文档，还深入理解了编译引擎的差异、字体处理的底层逻辑，以及像 \texttt{\textbackslash end} 这样的控制序列在参数中使用的禁忌。以后再遇到类似错误，就不需要死记硬背，而是能够\textbf{分析}——\textbf{切换引擎}、\textbf{补装材料}或\textbf{检查特殊字符}，从而游刃有余。

\begin{flushright}
  \emph{—— 编译愉快，学习顺利！}
\end{flushright}

\end{document}
